% GENERATED, edit the French unit or the translation table
\mxcorrections{4}\label{C04-CORR}
\textbf{Teacher use.} The answers below concern the tasks not worked in the lesson. First have students state their method; give graded help before the solution. The durations are estimates to be calibrated in class.
\subsection*{Diagnostic and remediation}
\begin{enumerate}
\item $\dfrac{8-2}{3-1}=3$. \textit{Action:} Represent the horizontal and vertical increments.
\item $4h+h^2$. \textit{Action:} Expand before simplifying.
\item $4+h$. \textit{Action:} Factor h out of the numerator before dividing.
\item $0$ and $\num{0,1}$. \textit{Action:} Check by squaring.
\item $3$ and $2$. \textit{Action:} Represent the distance to 0.
\item $y=2x+3$. \textit{Action:} Substitute the coordinates to find the y-intercept.
\end{enumerate}
\subsection*{Investigation activity}
The slopes are $2+h$; limit 2. The line passes through $\pt{1;1}$, so $y=2x-1$. The drawing motivates, the calculation justifies.
\subsection*{Immediate applications}
\begin{enumerate}
\item $f(2+h)-f(2)=7h+h^2$; rate $7+h$ for $h\ne0$.
\item The rate is $4+3h$, so the derivative at 1 is 4.
\item $y=5-2(x-3)=-2x+11$; approximation $5-2(\num{0,04})=\num{4,92}$.
\item The rate at 2 is $\dfrac{|h|}{h}$, with two opposite limiting values. For the square root at 0, the rate $\dfrac{1}{\sqrt h}$ for $h>0$ exceeds every bound as h gets closer to 0: no real derivative.
\end{enumerate}
\subsection*{Exercises: method and results}
\paragraph{Exercise 4.1 -- Two secants.} Slopes 3 and \num{2,1}; derivative at 1 is 2. The second is closer to the slope of the tangent; this is an exact calculation of a rate.
\paragraph{Exercise 4.2 -- A tangent.} Rate $4a+2h=-4+2h$, so $f\prime(-1)=-4$. $f(-1)=3$; tangent $y=3-4(x+1)=-4x-1$.
\paragraph{Exercise 4.3 -- Marginal cost.} $C(10)=400$, $C(11)=441$, increase 41 for one unit. The rate around 10 is $40+h$, so derivative 40: an approximation of the increase for a small change, not an equality for a whole unit.
\paragraph{Exercise 4.4 -- Reading a tangent.} On the figure in the lesson: points (1;0) and (2;4), slope 4. In the new situation, $f(2)=3$, slope $-\dfrac{1}{2}$. Approximation $3-\num{0,5}\times\num{0,1}=\num{2,95}$. The y-coordinate of B is not necessarily $f(4)$.
\paragraph{Exercise 4.5 -- Checking an approximation.} $1+2(-\num{0,02})=\num{0,96}$. Exact $\num{0,9604}$; exact minus approximate $\num{0,0004}>0$. The omitted term is $h^2$: approximation from below.
\paragraph{Exercise 4.6 -- Tangent and intersection.} The tangent has equation $y=1+3(x-1)=3x-2$. Expanding the product indeed gives $x^3-3x+2$. The x-coordinates of the intersections are therefore 1 and $-2$, by the zero-product rule. The points are $\pt{1;1}$ and $\pt{-2;-8}$. There are two common points: that is enough to refute the general claim.
\paragraph{Exercise 4.7 -- Numerical secants.} List of rates $2+h$: \num{2,1}; \num{2,01}; \num{2,001} and \num{1,9}; \num{1,99}; \num{1,999}. The algebraic calculation gives the real number 2. For very small h, floating-point calculation can lose precision.
\paragraph{Exercise 4.8 -- Delayed review.} For $h\ne0,-1$, the rate is $-\dfrac{1}{1+h}$; it tends to $-1$. Approximation $1-\num{0,02}=\num{0,98}$. The exact value $\dfrac{50}{51}$ allows a check.
\paragraph{CQFD, exercise 17 p.~127, all questions.} The x-coordinates of A, B and C are 0, 1 and 2. The tangents give slopes 2, 0 and $-2$ respectively. Thus $f\prime(0)=2$, $f\prime(1)=0$, $f\prime(2)=-2$. The units of the axes must be read from the grid.
\paragraph{CQFD, exercise 19 p.~127, questions 1 and 2.} The point A reads $\pt{1;4}$, so $f(1)=4$. The tangent passes through $\pt{0;1}$ and $\pt{1;4}$: slope 3. Thus $f\prime(1)=3$ and the tangent has equation $y=4+3(x-1)=3x+1$.
\subsection*{MCQ: answers and reasons}
\textbf{MCQ 1.}
\begin{enumerate}
\item B. $\dfrac{4-1}{2-1}=3$.
\item A. $4+2(x-1)$.
\item D. It is a value at a point.
\item C. Rate $6+h$.
\item B. Tangent affine function.
\end{enumerate}
\textbf{MCQ 2.}
\begin{enumerate}
\item C. Rate $\dfrac{1}{\sqrt h}$ not finite.
\item A. $|h|=-h$.
\item D. Constant y-coordinate.
\item B. A change in position divided by a change in time is expressed in $\mathrm{m\,s^{-1}}$.
\item C. Exact equality for an affine function, approximation in general.
\end{enumerate}
\subsection*{True or false}
\begin{enumerate}
\item False: $x^3$ and the tangent $3x-2$ at 1 also meet at $-2$.
\item False: the rates on either side are $-1$ and 1.
\item True: the rate is exactly m for every nonzero h.
\end{enumerate}
\paragraph{Differentiated teaching.} For a technical block, offer only the first step (domain, factor or reference sample space), then withdraw the help. For an independent student, ask for a second method and the treatment of a limiting case. Orally, allow an initial reformulation in the language the student has mastered, then require a precise mathematical sentence in French.
